%% ======================================================================
%%  A supplement, without proofs  (pp. 94-96)
%% ======================================================================
\clearpage
\folio{94}%
\global\let\currentseccode\empty
\global\let\pagefirstsec\empty
\setcounter{equation}{0}%
\renewcommand{\chaptertitle}{}%
\def\currentsectitle{}%
\thispagestyle{emptyhead}%
\phantomsection
\pdfbookmark[0]{A supplement, without proofs}{fm:supplement}%
\label{fm:supplement}%
\addtocontents{toc}{\string\fml{A SUPPLEMENT, WITHOUT PROOFS}{\thepage}{fm:supplement}}%
\begin{center}{\Large A supplement, without proofs}\end{center}%
\vspace{1.2\baselineskip}%
\par

This appendix presents, without proofs, some definitions and facts which
complete and generalize certain ideas discussed in the text. The notions to
be introduced here are not much deeper or more difficult than the
corresponding material in the text. The author has kept a detailed
discussion of these matters out of the text in order to keep the book short.
Besides, the reader may find these more general results in numerous
reference sources (see, for example, Chapter V of \hyperref[bib:9]{[9]}, which may be read
immediately after this book).

It has not been the author's intention to summarize here, without proofs,
developments in number theory which lie beyond the scope of this book
(adèles, class fields, zeta functions and $L$-series, the arithmetic of
simple algebras, analytic number theory, quadratic forms, etc.). He has
preferred rather to refer the reader to the works cited in the first part of
the bibliography (the unnumbered titles).

\divhead{Transitivity formulas}

Given three fields, $K\subset L\subset M$, each extension of finite degree
over the preceding, we have ``trace'' mappings
\[
\operatorname{Tr}_{L/K}:L\to K,\quad
\operatorname{Tr}_{M/L}:M\to L,\quad
\operatorname{Tr}_{M/K}:M\to K,
\]
and, analogously, ``norm'' mappings (\hyperref[sec:2.6]{Chapter II, §6}). For any $x\in M$, we
have the following relations
\begin{equation}
\left\{
\begin{aligned}
\operatorname{Tr}_{M/K}(x)&=\operatorname{Tr}_{L/K}
(\operatorname{Tr}_{M/L}(x))\\
\mathrm{N}_{M/K}(x)&=\mathrm{N}_{L/K}(\mathrm{N}_{M/L}(x)).
\end{aligned}
\right.
\label{eq:Supp.1}
\end{equation}

\divhead{The relative norm of an ideal}

Given two number fields, $K\subset K'$, and an ideal (integral or
fractional) $\mathfrak{a}'$ of $K'$, the ideal of $K$ generated by the
elements of $K$ of the form $\mathrm{N}_{K'/K}(x)$ ($x\in\mathfrak{a}'$) is
called the \emph{relative norm} of $\mathfrak{a}'$. It is denoted
$\mathrm{N}_{K'/K}(\mathfrak{a}')$ (or $\mathrm{N}(\mathfrak{a}')$). If
$\mathfrak{a}'$ is a principal ideal $(\alpha')$, then
\begin{equation}
\mathrm{N}_{K'/K}(\alpha')=(\mathrm{N}_{K'/K}(\alpha')).
\label{eq:Supp.2}
\end{equation}

If $K=\mathbf{Q}$, this definition of norm coincides with that presented in
\hyperref[sec:3.5]{Chapter III, §5}. If $\mathfrak{a}'$ is an integral ideal of $K'$ and if $A'$
is the ring of integers of $K'$, then
\begin{equation}
\mathrm{N}_{K'/\mathbf{Q}}(\mathfrak{a}')
=\operatorname{card}(A'/\mathfrak{a}')\mathbf{Z}.
\label{eq:Supp.3}
\end{equation}

Let us return to the general case. If $\mathfrak{a}$ is an ideal of $K$ and
if $n=[K':K]$, then
\begin{equation}
\mathrm{N}_{K'/K}(A'\mathfrak{a})=\mathfrak{a}^{n}.
\label{eq:Supp.4}
\end{equation}

For $\mathfrak{a}'$ and $\mathfrak{b}'$ ideals of $K'$ there is the
multiplicativity formula
\begin{equation}
\mathrm{N}_{K'/K}(\mathfrak{a}'\mathfrak{b}')
=\mathrm{N}_{K'/K}(\mathfrak{a}')\mathrm{N}_{K'/K}(\mathfrak{b}').
\label{eq:Supp.5}
\end{equation}

Finally, if $\mathfrak{p}'$ is a prime ideal of $K'$, if
$\mathfrak{p}=\mathfrak{p}'\cap K$, and if $f$ is the residual degree of
$\mathfrak{p}'$ over $K$, then
\begin{equation}
\mathrm{N}_{K'/K}(\mathfrak{p}')=\mathfrak{p}^{f}.
\label{eq:Supp.6}
\end{equation}

For three number fields $K\subset K'\subset K''$ there is the following
transitivity formula, in which $\mathfrak{a}''$ denotes an ideal of $K''$:
\begin{equation}
\mathrm{N}_{K''/K}(\mathfrak{a}'')
=\mathrm{N}_{K'/K}\bigl(\mathrm{N}_{K''/K'}
(\mathfrak{a}'')\bigr).
\label{eq:Supp.7}
\end{equation}

In the same context, in terms of the relative norm of an ideal one may
obtain a transitivity formula for discriminants (where
$\mathcal{D}_{K'/K}$ denotes the discriminant of $K'$ over $K$; cf.\
\hyperref[sec:5.3]{Chapter V, §3}, \hyperref[def:V.3.1]{Definition 1}):
\begin{equation}
\mathcal{D}_{K''/K}
=\mathrm{N}_{K'/K}(\mathcal{D}_{K''/K'})
\mathcal{D}_{K'/K}^{[K'':K']}.
\label{eq:Supp.8}
\end{equation}

The above generalizes to the case of a Dedekind ring $A$ and its integral
closure $A'$ in a finite extension of the field of fractions of $A$.

\divhead{The different}

The following is valid for a Dedekind ring $A$ and its integral closure in a
finite extension of its field of fractions. For simplicity we limit our
discussion to the number field case.

Let $K$ and $K'$ be fields, $K\subset K'$. Let $A$ and $A'$ be the respective
rings of integers in $K$ and $K'$. One says that a maximal ideal
$\mathfrak{p}'$ of $A'$ is \emph{ramified over} $A$ (or over $K$) if its
index of ramification over $A$ is larger than $1$. In this case one says
that the maximal ideal $\mathfrak{p}=\mathfrak{p}'\cap A$ of $A$
\emph{ramifies in} $A'$ (\hyperref[sec:5.3]{Chapter V, §3}). It follows easily from \hyperref[sec:5.3]{Chapter V,
§3}, \hyperref[thm:V.3.1]{Theorem 1} that the maximal ideals of $A'$ which are ramified over $A$
are finite in number. We are going to describe an ideal
$\mathfrak{d}_{K'/K}$ of $A'$, called the ``different'' of $K'$ over $K$,
which will have the property that the maximal ideals ramified over $A$ are
precisely those which contain $\mathfrak{d}_{K'/K}$ (note the analogy with
\hyperref[sec:5.3]{Chapter V, §3}, \hyperref[thm:V.3.1]{Theorem 1}).

First, one proves that the set of all $x\in K'$ such that
\begin{equation}
\operatorname{Tr}_{K'/K}(xA')\subset A
\label{eq:Supp.9}
\end{equation}
is a fractional ideal $\mathfrak{C}$ of $A'$; it is called the
\emph{codifferent} of $K'$ over $K$. By definition the different
$\mathfrak{d}_{K'/K}$ is the inverse ideal $\mathfrak{C}^{-1}$. The
different is a non-zero \emph{integral} ideal of $A'$.

It can be shown that the different is generated by elements of the form
$F'(x)$, where $x$ runs through $A'$ and where $F$ denotes the minimal
polynomial of $x$ over $K$ and where $F'$ is the derivative of $F$. In
particular, if $A'$ is of the form $A[y]$ (which isn't always the case) and
if $G$ is the minimal polynomial for $y$ over $K$, then the different
$\mathfrak{d}_{K'/K}$ is the principal ideal of $A'$ generated by $G'(y)$.

The non-zero prime ideals of $A'$ which are ramified over $A$ are those
which contain $\mathfrak{d}_{K'/K}$. More precisely, let
\begin{equation}
\mathfrak{d}_{K'/K}=\prod_{i}\mathfrak{p}_{i}^{m_{i}}
\quad(m_{i}>0)
\label{eq:Supp.10}
\end{equation}
be the decomposition of the different into a product of powers of prime
ideals, and let $e_{i}$ be the index of ramification of
$\mathfrak{p}_{i}'$ over $A$. Then the non-zero prime ideals of $A'$ which
are ramified over $A$ are the $\mathfrak{p}_{i}'$ and
$m_{i}\geq e_{i}-1$ for all $i$. Furthermore, $m_{i}=e_{i}-1$ if and only
if $e_{i}$ is prime to the characteristic of the residual field
$A'/\mathfrak{p}_{i}'$.

The different $\mathfrak{d}_{K'/K}$ (ideal of $A'$) and the discriminant
$\mathcal{D}_{K'/K}$ (ideal of $A$) are connected by the following relation
\begin{equation}
\mathcal{D}_{K'/K}=\mathrm{N}_{K'/K}(\mathfrak{d}_{K'/K})
\label{eq:Supp.11}
\end{equation}
(cf.\ \hyperref[sec:2.7]{Chapter II, §7}, formula \eqref{eq:II.7.6}). Thus the different gives
more precise information concerning the ramification of primes in $A$ and
$A'$ than does the discriminant.

Finally, for three number fields, $K\subset K'\subset K''$, there is the
following transitivity relation for the differents
\begin{equation}
\mathfrak{d}_{K''/K}=\mathfrak{d}_{K''/K'}
\mathfrak{d}_{K'/K}.
\label{eq:Supp.12}
\end{equation}
